\lysh{16810_SOLUTION}{1}
\text{ΛΥΣΗ}

% TODO: TikZ σχήμα (μην το κατασκευάσεις) - The image shows a coordinate system with points A(1,1), B(5,2), Γ(0,-2) and Δ(8,0) plotted, forming a quadrilateral.

\text{α) Τοποθετούμε τα σημεία στο επίπεδο όπως φαίνεται στο παραπάνω σχήμα. Για να είναι}
\text{το τετράπλευρο ΑΒΓΔ τραπέζιο αρκεί να αποδείξουμε ότι οι πλευρές ΑΒ και ΓΔ είναι}
\text{παράλληλες και οι πλευρές ΑΓ και ΒΔ τέμνονται.}

\text{Είναι } \lambda_{\text{AB}} = \frac{2-1}{5-1} = \frac{1}{4} \text{ και } \lambda_{\text{ΓΔ}} = \frac{0-(-2)}{8-0} = \frac{0+2}{8-0} = \frac{2}{8} = \frac{1}{4} \text{. Άρα AB // ΓΔ.}

\text{Επίσης, } \lambda_{\text{AΓ}} = \frac{-2-1}{0-1} = \frac{-3}{-1} = 3 \text{ και } \lambda_{\text{BΔ}} = \frac{0-2}{8-1} = \frac{-2}{7} = -\frac{2}{7} \text{. Άρα } \lambda_{\text{AΓ}} \neq \lambda_{\text{BΔ}} \text{ και οι ευθείες ΑΓ και ΒΔ δεν είναι}
\text{παράλληλες.}

\text{β) Για το εμβαδόν του τραπεζίου ΑΒΓΔ έχουμε: } (\text{ΑΒΓΔ}) = (\text{ΑΒΔ}) + (\text{ΑΓΔ}) \ (1).

\text{Για τα εμβαδά των τριγώνων ΑΒΔ και ΑΓΔ υπολογίζουμε πρώτα τα διανύσματα } \vec{\text{ΑΓ}}, \vec{\text{ΑΔ}} \text{ και}
$\vec{\text{AB}}$ \text{ και έχουμε:}
\[
\vec{\text{AΓ}} = (0-1, -2-1) = (-1, -3)
\]
\[
\vec{\text{AΔ}} = (8-1, 0-1) = (7, -1)
\]
\[
\vec{\text{AB}} = (5-1, 2-1) = (4, 1)
\]
\[
\det(\vec{\text{AΓ}}, \vec{\text{AΔ}}) = \begin{vmatrix} -1 & -3 \\ 7 & -1 \end{vmatrix} = (-1)(-1) - (-3)(7) = 1 + 21 = 22
\]
\[
\det(\vec{\text{AB}}, \vec{\text{AΔ}}) = \begin{vmatrix} 4 & 1 \\ 7 & -1 \end{vmatrix} = (4)(-1) - (1)(7) = -4 - 7 = -11
\]
\[
(\text{ΑΒΔ}) = \frac{1}{2} |\det(\vec{\text{AB}}, \vec{\text{AΔ}})| = \frac{1}{2} |-11| = \frac{11}{2} \text{ τ.μ.}
\]
\[
(\text{ΑΓΔ}) = \frac{1}{2} |\det(\vec{\text{AΓ}}, \vec{\text{AΔ}})| = \frac{1}{2} |22| = 11 \text{ τ.μ.}
\]
\text{Οπότε η σχέση (1) γίνεται: } $(\text{ΑΒΓΔ}) = \frac{11}{2} + 11 = \frac{11}{2} + \frac{22}{2} = \frac{33}{2} \text{ τ.μ.}$
\end{lysh}